\documentclass[10pt,a4paper]{article}

\usepackage[margin=0.72in]{geometry}
\usepackage[T1]{fontenc}
\usepackage[utf8]{inputenc}
\usepackage{booktabs}
\usepackage{graphicx}
\usepackage{xcolor}
\usepackage{hyperref}
\usepackage{enumitem}
\usepackage{array}

\hypersetup{hidelinks}
\setlist{nosep,leftmargin=*}
\setlength{\parindent}{0pt}
\setlength{\parskip}{5pt}
\renewcommand{\arraystretch}{1.08}
\newcommand{\demo}[1]{\textcolor{blue!65!black}{#1\textsuperscript{*}}}
\newcommand{\pending}{\textcolor{gray}{--}}

\title{\vspace{-1.4cm}LLM-DNA Robustness Experiment\\
\large Five-Page Progress Brief}
\author{}
\date{Updated July 30, 2026}

\begin{document}
\maketitle
\vspace{-0.8cm}

\section{What changed and what matters}

\begin{quote}
\textbf{Main finding.} The current 84-model exploratory evidence does not
support the simple claim that higher temperature always reduces model-detection
performance. Instead, temperature and top-$p$ interact: increasing $T$ from
0.2 to 0.3 improved performance at $p=0.8$, made little detectable difference
at $p=0.9$, and substantially reduced performance at $p=1.0$.
\end{quote}

\textbf{Changes from the previous report}

\begin{itemize}
  \item The primary cohort is now the fixed validated-100 cohort.
  \item The plan was reduced from four to two generation rounds per condition:
        repeat 1 is the reference/training split and repeat 2 is held out.
  \item The workload is now 13 conditions $\times$ 2 rounds $\times$ 100
        models $=2{,}600$ successful artifacts, rather than 5,200.
  \item The report now includes the live grid, exploratory values in the main
        result tables, direct answers to the research questions, and a dated
        completion forecast.
\end{itemize}

\textbf{Current status (July 30, 18:16 Singapore time).}
The live audit contains \textbf{1,352/2,600 successful artifacts (52.0\%)},
leaving 1,248. A four-GPU backfill on GPUs 0, 1, 2, and 4 is active and has
added 39 successes since its July 30 restart. No primary accuracy claim will
be made until both rounds are complete on a common cohort.

\begin{table}[h]
\centering
\small
\caption{Current answers to the research questions.}
\begin{tabular}{p{0.08\linewidth}p{0.58\linewidth}p{0.25\linewidth}}
\toprule
RQ & Current answer & Confidence \\
\midrule
RQ1 & Temperature cannot be interpreted independently of top-$p$. &
Limited 84-model evidence \\
RQ2 & Linear and RBF SVMs improve equal-cell Top-1 over cosine on identical
queries. & Paired exploratory CIs are positive \\
RQ3 & The proposed distributional advantage is not yet tested fairly on the
84-model cohort. & Inconclusive \\
RQ4 & Linear SVM has the best mixed accuracy; RBF has the best worst-cell
accuracy at $\tau=0.3$. & Limited 84-model evidence \\
\bottomrule
\end{tabular}
\end{table}

\clearpage
\section{Important exploratory findings}

The following results use the same 84 availability-selected models, repeat 1
for pooled reference/training construction, and repeat 2 for held-out testing.
They are pipeline-scale evidence, not final validated-100 results.

\begin{table}[h]
\centering
\small
\caption{Held-out exact-model Top-1. Asterisks denote exploratory results.}
\begin{tabular}{ccrrrr}
\toprule
$T$ & Top-$p$ & Cosine & Linear SVM & RBF SVM & Proposed \\
\midrule
0.2 & 0.8 & \demo{0.643} & \demo{0.821} & \demo{0.702} & \pending \\
0.2 & 0.9 & \demo{0.726} & \demo{0.845} & \demo{0.762} & \pending \\
0.2 & 1.0 & \demo{0.595} & \demo{0.643} & \demo{0.643} & \pending \\
0.3 & 0.8 & \demo{0.774} & \demo{0.869} & \demo{0.857} & \pending \\
0.3 & 0.9 & \demo{0.690} & \demo{0.786} & \demo{0.762} & \pending \\
0.3 & 1.0 & \demo{0.298} & \demo{0.369} & \demo{0.405} & \pending \\
\bottomrule
\end{tabular}
\end{table}

\textbf{Temperature--top-$p$ interaction.}
For cosine retrieval, increasing $T$ from 0.2 to 0.3 changed Top-1 by
$+0.131$ at $p=0.8$ (95\% CI 0.024 to 0.238), $-0.036$ at $p=0.9$
($-0.167$ to 0.083), and $-0.298$ at $p=1.0$ ($-0.417$ to $-0.179$).
The two SVMs show the same qualitative pattern. The current data therefore
support an interaction at limited scale, but not a general monotone
temperature conclusion.

\begin{table}[h]
\centering
\small
\caption{Temperature-agnostic bounded deployment.}
\begin{tabular}{clrr}
\toprule
$\tau$ & Method & Mixed Top-1 & Worst cell \\
\midrule
0.2 & Cosine & 0.596 & 0.562 \\
    & Linear SVM & \textbf{0.682} & \textbf{0.652} \\
    & RBF SVM & 0.629 & 0.584 \\
\addlinespace
0.3 & Cosine & 0.621 & 0.298 \\
    & Linear SVM & \textbf{0.722} & 0.369 \\
    & RBF SVM & 0.688 & \textbf{0.405} \\
\bottomrule
\end{tabular}
\end{table}

On identical queries at $\tau=0.3$, linear SVM improved over cosine by 0.101
(paired 95\% CI 0.058 to 0.147), and RBF SVM improved by 0.067
(0.026 to 0.111). Thus classifier mitigation is the strongest finding
currently supported, although only at exploratory scale.

\clearpage
\section{Live generation grid}

Each planned setting requires 200 successes: 100 models in each of two rounds.
A dash denotes an unplanned deterministic duplicate.

\begin{table}[h]
\centering
\small
\caption{Reusable successes as of July 30, 18:16.}
\begin{tabular}{crrr}
\toprule
Temperature & $p=0.8$ & $p=0.9$ & $p=1.0$ \\
\midrule
$0$   & -- & -- & $42/200$ \\
$0.2$ & $166/200$ & $171/200$ & $162/200$ \\
$0.3$ & $170/200$ & $168/200$ & $161/200$ \\
$0.5$ & $79/200$  & $85/200$  & $73/200$ \\
$0.7$ & $1/200$   & $74/200$  & $0/200$ \\
\bottomrule
\end{tabular}
\end{table}

\begin{table}[h]
\centering
\small
\caption{Progress summarized by temperature.}
\begin{tabular}{lrrr}
\toprule
Condition & Required & Success & Remaining \\
\midrule
$T=0$ & 200 & 42 & 158 \\
$T=0.2$, all $p$ & 600 & 499 & 101 \\
$T=0.3$, all $p$ & 600 & 499 & 101 \\
$T=0.5$, all $p$ & 600 & 237 & 363 \\
$T=0.7$, all $p$ & 600 & 75 & 525 \\
\midrule
Total & 2,600 & 1,352 & 1,248 \\
\bottomrule
\end{tabular}
\end{table}

\textbf{Execution priority.}
Resume the deterministic control first, followed by the large $T=0.7$ gaps,
then $T=0.5$, and finally the smaller $T=0.2/0.3$ backfills. Successful
historical artifacts are preserved; only missing or failed model--cell--round
keys are queued.

\textbf{What is and is not complete.}
All 13 settings have been started, but none yet contains 100 successful models
in both rounds. The filled accuracy values on the previous page come from a
separate 84-model demonstration and must not be pooled with the primary
validated-100 results.

\textbf{Current execution.}
The four-GPU backfill remains active. At the 18:20 host-side check, four model
workers were still processing the deterministic $r1$ cell. The status journal
preserves the 42 completed deterministic artifacts, so later restarts will not
regenerate them.

\clearpage
\section{Revised design and the proposed method}

\textbf{Primary design}

\begin{itemize}
  \item Cohort: fixed validated 100 Hugging Face models.
  \item Prompts: the same 100 classical-Chinese prompts for every model.
  \item Conditions: $T=0$ once, plus
        $T\in\{0.2,0.3,0.5,0.7\}\times p\in\{0.8,0.9,1.0\}$.
  \item Split: repeat 1 is reference/training; repeat 2 is held-out query.
  \item Primary metric: exact-model Top-1; Top-3, Top-5, and MRR are secondary.
  \item Deployment view: decoding parameters are hidden, with only
        $T\leq\tau$ known.
\end{itemize}

\textbf{Why the Proposed column is still blank.}
The same 84 models and held-out queries can be evaluated with exact RBF-MMD at
$K=1$. However, one reference and one query response per prompt do not test a
multi-sample distribution: at $K=1$, MMD is effectively a kernel point
distance. In addition, the existing eight-model pilot used character
n-gram hashing rather than the response encoder used for the main DNA table.
Its saturated 1.000 accuracy is therefore neither difficult nor directly
comparable.

\begin{table}[h]
\centering
\small
\caption{Data required for a disjoint, matched-budget distributional test.}
\begin{tabular}{lcc}
\toprule
Budget & Reference samples per prompt & Query samples per prompt \\
\midrule
$K=1$ & 1 & 1 \\
$K=2$ & 2 & 2 \\
$K=4$ & 4 & 4 \\
$K=8$ & 8 & 8 \\
\bottomrule
\end{tabular}
\end{table}

\textbf{Defensible next step for the proposed method.}
First run an explicitly labelled ``RBF-MMD, $K=1$'' analysis on the exact
84-model cohort using the same encoder, prompts, repeat split, and queries.
This can check the paired pipeline without new generation. A claim about the
proposed multi-sample advantage requires additional independent responses and
matched $K\geq2$ cosine/mean baselines; it should not be inferred from the
current table.

\textbf{Main limitations}

\begin{itemize}
  \item The exploratory cohort was selected by artifact availability.
  \item Only $T=0.2$ and 0.3 currently have comparable held-out results.
  \item Two rounds quantify model-level uncertainty, not variability across
        independent generation rounds.
  \item The shared $T=0.5$ seed-42 runs require a duplicate/independence audit.
\end{itemize}

\clearpage
\section{Schedule, decisions, and next report}

The latest completed four-GPU batch produced 73 successes in 12.43 hours
(5.87 successes/hour). At that rate, 1,248 remaining artifacts require about
213 four-GPU hours, or 8.9 days of continuous generation, before retry and
audit allowances.

\begin{table}[h]
\centering
\small
\caption{Conditional completion forecast.}
\begin{tabular}{p{0.32\linewidth}p{0.20\linewidth}p{0.38\linewidth}}
\toprule
Milestone & Expected & Completion criterion \\
\midrule
Missing/failed backfill & In progress & Four workers active on GPUs
0, 1, 2, and 4 since July 30 \\
Generation and retries & Aug.\ 11--12 & 2,600 admissible successes \\
Independence/provenance audit & Aug.\ 12 & Split passes seed and duplicate checks \\
Primary evaluation and figures & Aug.\ 13--14 & Final common-cohort outputs \\
Revised primary report & Aug.\ 14 & Exploratory values replaced by primary results \\
\bottomrule
\end{tabular}
\end{table}

\textbf{Schedule assessment.}
August 14 remains feasible only if the active four-GPU backfill continues at
approximately the observed throughput. Every
additional idle day moves the forecast by at least one day. Persistent
$T=0.7$ failures or a distinct-seed $T=0.5$ backfill would also extend it.

\textbf{Decisions for the next report}

\begin{enumerate}
  \item Report temperature effects separately by top-$p$; do not present only
        a marginal temperature coefficient.
  \item Treat classifier mitigation as exploratory until the validated-100
        common cohort is complete.
  \item Add the proposed method only under a clearly stated sample budget and
        matched encoder/query split.
  \item Keep the main report short; move protocol commands, manifests, and
        extended validity checks to an appendix or repository files.
\end{enumerate}

\textbf{Reproducibility pointers.}
The fixed cohort is
\path{configs/rand_chinese_validated100_latest.jsonl}; the live machine-readable
grid is \path{reports/primary_robustness_run_grid.csv}; exploratory predictions
are under \path{out/small_scale_section_6_11}; and the backfill log is
\path{logs/primary_grid_resume_20260730-1235_master.log}.

\vfill
\begin{center}
\textit{Primary conclusions remain pending until the validated-100 grid and
held-out split are complete.}
\end{center}

\end{document}
